% Typografie-Atelier · DIN A4 · Eingaben sind Text, keine TeX-Befehle.
% Engine: pdfLaTeX
\documentclass[11pt,a4paper,ngerman]{scrlttr2}
\usepackage[T1]{fontenc}
\usepackage[utf8]{inputenc}
\usepackage{tgpagella}
\usepackage{tgheros}
\usepackage[ngerman]{babel}
\usepackage{microtype}
\setkomafont{subject}{\normalfont\bfseries}
\setkomavar{fromname}{Beispielname}
\setkomavar{date}{9. Oktober 2026}
\setkomavar{subject}{Ein Raum für Gedanken}
\begin{document}
\begin{letter}{Beispielschule\\
Musterstraße 1\\
12345 Musterstadt}
\opening{Sehr geehrte Damen und Herren,}
Ein guter Text lädt zum Weiterdenken ein. Die Gestaltung ordnet Gedanken, ohne sie zu überlagern.

Schriftwahl, Zeilenlänge und Durchschuss bilden ein System. Ein Preis von 10 € oder 50 \% wird hier als normaler Text behandelt.
\closing{Mit freundlichen Grüßen}
\end{letter}
\end{document}
